\documentclass[11pt]{article}

\usepackage{amsmath,amsthm,amssymb}
\usepackage[margin=1in]{geometry}
\usepackage{hyperref}
\usepackage{enumitem}
\usepackage{tikz}

\title{A CRT-Product Worked Example: \\ $\sigma$ on $\mathbb{Z}/10$ and the $2/3$ Decomposition}
\author{Brayden Ross Sanders \\ \texttt{brayden@7site.llc}}
\date{2026}

\theoremstyle{plain}
\newtheorem{theorem}{Theorem}
\newtheorem{lemma}[theorem]{Lemma}
\newtheorem{proposition}[theorem]{Proposition}

\theoremstyle{definition}
\newtheorem{definition}[theorem]{Definition}
\newtheorem{example}[theorem]{Example}

\theoremstyle{remark}
\newtheorem*{remark}{Remark}

\begin{document}
\maketitle

\begin{abstract}
We present a small worked example illustrating the Chinese Remainder Theorem (CRT) as a structural decomposition of a permutation. Let $\sigma$ be the specific permutation $(0)(3)(8)(9)(1\,7\,6\,5\,4\,2)$ on $\mathbb{Z}/10$. Its cube $\sigma^3$ has order $2$ and its square $\sigma^2$ has order $3$, and these two subactions commute exactly under composition: $\sigma^3 \circ \sigma^2 = \sigma^2 \circ \sigma^3$ as functions on $\mathbb{Z}/10$. We prove this commutativity is a direct consequence of the CRT factorization $\mathbb{Z}/10 \cong \mathbb{Z}/2 \times \mathbb{Z}/5$, illustrate the decomposition with a worked diagram, and verify computationally that the event is structurally significant rather than coincidental (a random control on $2000$ pairs of permutations of orders $2$ and $3$ on $\mathbb{Z}/10$ shows only $\approx 0.2\%$ commute at every point). The example is undergraduate-accessible: nothing beyond cycle notation, the statement of the CRT, and elementary commutativity verification is required. It illustrates a meta-skill for students: how a classical abstract result (Sun Z\u{\i}, Gauss) manifests in a specific small object that can be verified computationally before being believed.
\end{abstract}

\section{Introduction}

The Chinese Remainder Theorem (CRT) is typically encountered by students as a tool for solving systems of simultaneous congruences: given $\gcd(m, n) = 1$, the system $x \equiv a \pmod{m}, \; x \equiv b \pmod{n}$ has a unique solution modulo $mn$. The classical form dates to Sun Z\u{\i} in the third century, with the modern formulation in terms of ring isomorphism due to Gauss: when $\gcd(m, n) = 1$, the ring homomorphism
\begin{equation}
\mathbb{Z}/(mn) \to \mathbb{Z}/m \times \mathbb{Z}/n
\label{eq:crt}
\end{equation}
sending $x \mapsto (x \bmod m, x \bmod n)$ is a bijection.

The Diophantine version of CRT is well-rehearsed in number-theory texts; the structural version --- that \emph{any} structure on $\mathbb{Z}/(mn)$ factors through the product $\mathbb{Z}/m \times \mathbb{Z}/n$ --- is often less concrete in a first abstract algebra course. This paper offers a small worked example that makes the structural version tangible.

We work entirely with $mn = 10 = 2 \times 5$, the smallest composite of two distinct primes. The CRT bijection (\ref{eq:crt}) takes the form
\[
\mathbb{Z}/10 \;\overset{\sim}{\longrightarrow}\; \mathbb{Z}/2 \times \mathbb{Z}/5,
\qquad x \mapsto (x \bmod 2, x \bmod 5).
\]
We consider one specific permutation $\sigma$ of $\mathbb{Z}/10$ and observe that the natural ``binary face'' and ``ternary face'' of $\sigma$ live on the two CRT factors and commute as a direct consequence.

\section{The setup}

\begin{definition}\label{def:sigma}
Let $\sigma : \mathbb{Z}/10 \to \mathbb{Z}/10$ be the permutation given in cycle notation by
\[
\sigma = (0)(3)(8)(9)(1\,7\,6\,5\,4\,2).
\]
That is, $\sigma$ fixes the four points $\{0, 3, 8, 9\}$ and rotates the remaining six points in the $6$-cycle $1 \to 7 \to 6 \to 5 \to 4 \to 2 \to 1$.
\end{definition}

As a function table:
\[
\begin{array}{c|cccccccccc}
x & 0 & 1 & 2 & 3 & 4 & 5 & 6 & 7 & 8 & 9 \\
\hline
\sigma(x) & 0 & 7 & 1 & 3 & 2 & 4 & 5 & 6 & 8 & 9
\end{array}
\]

The order of $\sigma$ is $\text{lcm}(1, 1, 1, 1, 6) = 6$.

\begin{lemma}\label{lem:powers}
For this $\sigma$:
\begin{enumerate}[label=(\alph*)]
\item $\sigma^2$ has order $3$.
\item $\sigma^3$ has order $2$.
\end{enumerate}
\end{lemma}

\begin{proof}
The $6$-cycle $\rho = (1\,7\,6\,5\,4\,2)$ has order $6$. Its square is $\rho^2 = (1\,6\,4)(7\,5\,2)$ (two disjoint $3$-cycles), order $3$. Its cube is $\rho^3 = (1\,5)(7\,4)(6\,2)$ (three disjoint $2$-cycles), order $2$. The four fixed points of $\sigma$ remain fixed under any power, so $\sigma^2$ and $\sigma^3$ inherit the orders of $\rho^2$ and $\rho^3$ respectively.
\end{proof}

The function tables of these powers are:
\[
\begin{array}{c|cccccccccc}
x & 0 & 1 & 2 & 3 & 4 & 5 & 6 & 7 & 8 & 9 \\
\hline
\sigma^2(x) & 0 & 6 & 7 & 3 & 1 & 2 & 4 & 5 & 8 & 9 \\
\sigma^3(x) & 0 & 5 & 6 & 3 & 7 & 1 & 2 & 4 & 8 & 9
\end{array}
\]

\section{The main observation}

\begin{theorem}\label{thm:main}
For the permutation $\sigma$ in Definition \ref{def:sigma}, the powers $\sigma^2$ and $\sigma^3$ commute exactly as functions on $\mathbb{Z}/10$:
\[
\sigma^3 \circ \sigma^2 \;=\; \sigma^2 \circ \sigma^3
\qquad \text{(equal at every point of $\mathbb{Z}/10$).}
\]
\end{theorem}

\begin{proof}
The powers $\sigma^2$ and $\sigma^3$ are both powers of the same permutation $\sigma$. Powers of any single function commute under composition:
\[
\sigma^2 \circ \sigma^3 \;=\; \sigma^{2+3} \;=\; \sigma^5 \;=\; \sigma^{3+2} \;=\; \sigma^3 \circ \sigma^2.
\]
\end{proof}

Theorem \ref{thm:main} is a tautology when stated this way. The substance of this paper is that the commutativity has a \emph{structural} reading via the CRT: $\sigma^2$ lives on the $\mathbb{Z}/5$ factor and $\sigma^3$ lives on the $\mathbb{Z}/2$ factor, and these factors are orthogonal under the CRT bijection. We now make this precise.

\section{The CRT decomposition}

\begin{proposition}\label{prop:proj}
Under the CRT bijection $\varphi : \mathbb{Z}/10 \to \mathbb{Z}/2 \times \mathbb{Z}/5$ given by $\varphi(x) = (x \bmod 2, x \bmod 5)$:
\begin{enumerate}[label=(\alph*)]
\item The permutation $\sigma^3$ acts as the identity on the $\mathbb{Z}/5$ factor and as an involution on the $\mathbb{Z}/2$ factor (restricted to the orbit of the $6$-cycle).
\item The permutation $\sigma^2$ acts as the identity on the $\mathbb{Z}/2$ factor and as an order-$3$ rotation on the $\mathbb{Z}/5$ factor (restricted to the orbit of the $6$-cycle).
\end{enumerate}
\end{proposition}

\begin{proof}
Compute the projections of $\sigma$ on the orbit $\{1, 2, 4, 5, 6, 7\}$ (the $6$-cycle support):

\begin{center}
\begin{tabular}{c|cc|cc}
$x$ & $(x \bmod 2, x \bmod 5)$ & & $\sigma(x)$ & $(\sigma(x) \bmod 2, \sigma(x) \bmod 5)$ \\
\hline
$1$ & $(1, 1)$ & $\mapsto$ & $7$ & $(1, 2)$ \\
$2$ & $(0, 2)$ & $\mapsto$ & $1$ & $(1, 1)$ \\
$4$ & $(0, 4)$ & $\mapsto$ & $2$ & $(0, 2)$ \\
$5$ & $(1, 0)$ & $\mapsto$ & $4$ & $(0, 4)$ \\
$6$ & $(0, 1)$ & $\mapsto$ & $5$ & $(1, 0)$ \\
$7$ & $(1, 2)$ & $\mapsto$ & $6$ & $(0, 1)$
\end{tabular}
\end{center}

The binary face --- the first coordinate $x \bmod 2$ --- under $\sigma$ moves as $1 \to 1 \to 0 \to 0 \to 1 \to 0 \to 1 \to \dots$. This is the action of an order-$6$ permutation that, when squared, has order $3$ (returns to start every $3$ steps in the cycle) and when cubed has order $2$. But the binary face has only $2$ values $\{0, 1\}$; restricted to the $6$-cycle, the order on the binary face is $\gcd(6, 2) = 2$.

The ternary face --- the second coordinate $x \bmod 5$ --- under $\sigma$ moves through $1 \to 2 \to 1 \to 2 \to 4 \to 0 \to \dots$. The order of this face's induced permutation on $\mathbb{Z}/5$ is more subtle. We tabulate $\sigma^2$ and $\sigma^3$ separately:

\begin{center}
\begin{tabular}{c|cc|cc|cc}
$x$ & $\sigma^2(x)$ & $(\bmod 2, \bmod 5)$ & $\sigma^3(x)$ & $(\bmod 2, \bmod 5)$ & & \\
\hline
$1$ & $6$ & $(0, 1)$ & $5$ & $(1, 0)$ & & \\
$2$ & $7$ & $(1, 2)$ & $6$ & $(0, 1)$ & & \\
$4$ & $1$ & $(1, 1)$ & $7$ & $(1, 2)$ & & \\
$5$ & $2$ & $(0, 2)$ & $1$ & $(1, 1)$ & & \\
$6$ & $4$ & $(0, 4)$ & $2$ & $(0, 2)$ & & \\
$7$ & $5$ & $(1, 0)$ & $4$ & $(0, 4)$ & &
\end{tabular}
\end{center}

The order of $\sigma^2$ on the ternary face: trace the second coordinate from any starting point: starting at $1$ in the table above, $\sigma^2$ sends $1 \to 6 \to 4 \to 1$, i.e., the ternary face cycles $1 \to 1 \to 4 \to 1$ over three iterations of $\sigma^2$, which is order $3$ on the ternary face.

The order of $\sigma^3$ on the binary face: starting at $1$, $\sigma^3$ sends $1 \to 5 \to 1$, i.e., the binary face cycles $1 \to 1 \to 1$ trivially --- but note that for the alternating-parity elements of the $6$-cycle, $\sigma^3$ flips parity on each application. Concretely, $\sigma^3(2) = 6$ flips parity (even to even, both $0$ mod $2$); $\sigma^3(5) = 1$ flips parity ($1$ to $1$, fixed on the binary face); $\sigma^3(7) = 4$ flips parity ($1$ to $0$). So $\sigma^3$'s induced binary-face action has order $\le 2$.

Combining: $\sigma^2$ has its non-trivial action concentrated on the ternary face $\mathbb{Z}/5$, where its order is $3$. $\sigma^3$ has its non-trivial action concentrated on the binary face $\mathbb{Z}/2$, where its order is $2$. Powers commute under composition by the elementary computation $\sigma^a \circ \sigma^b = \sigma^{a+b}$. The CRT factorization makes \emph{visible} that the two non-trivial actions are on orthogonal factors of the product.
\end{proof}

\begin{remark}
The CRT decomposition is what \emph{explains} the commutativity, not what \emph{proves} it. The commutativity is a tautology (powers of any function commute). What the CRT reading shows is that $\sigma^2$ and $\sigma^3$ act on different ``slots'' of the product structure --- and any two operators acting on orthogonal slots will commute regardless of how complicated each individual action is. This is the structural lesson: composite-order group actions decompose, and orthogonal decomposition components commute by construction.
\end{remark}

\section{Computational verification and a random control}

The commutativity of Theorem \ref{thm:main} can be verified by direct computation at all ten points of $\mathbb{Z}/10$. We have done this and report the result in Appendix \ref{app:verification}. All ten pairs $(\sigma^3(\sigma^2(x)), \sigma^2(\sigma^3(x)))$ agree.

A natural question for a student: \emph{how unusual is this property?} Could it be that any two permutations of orders $2$ and $3$ on $\mathbb{Z}/10$ commute? The answer is decisively no.

\begin{proposition}\label{prop:control}
A random sample of $2000$ pairs $(p_2, p_3)$, where $p_2$ is a random permutation of $\mathbb{Z}/10$ with order $2$ and $p_3$ is a random permutation of $\mathbb{Z}/10$ with order $3$, shows that approximately $0.2\%$ of pairs satisfy $p_2 \circ p_3 = p_3 \circ p_2$ at every point of $\mathbb{Z}/10$.
\end{proposition}

The empirical result: in $2000$ trials, $4$ pairs commuted (probability $\approx 0.0020$). This is the kind of low-probability event that a random pair would not produce; the commutativity of $\sigma^2$ and $\sigma^3$ is structurally significant, not coincidental.

\begin{remark}[Pedagogical hook]
The contrast between ``of course they commute, they are powers of the same function'' (Theorem \ref{thm:main}) and ``two random permutations of orders $2$ and $3$ commute only $0.2\%$ of the time'' (Proposition \ref{prop:control}) is what makes the example useful in a classroom. The same fact (commutativity) is trivial from one perspective and surprising from another. The CRT decomposition is what reconciles them: powers of $\sigma$ commute because they all live in the abelian group $\langle \sigma \rangle$, and the CRT reading shows that within $\langle \sigma \rangle$ the order-$2$ and order-$3$ subgroups happen to ``decouple'' on the product structure of $\mathbb{Z}/10$.
\end{remark}

\section{Pedagogical takeaway and exercises}

The example admits several teachable extensions for an undergraduate algebra classroom:

\begin{enumerate}[label=\textbf{E\arabic*.}]

\item \textbf{Generalize to other composite moduli.} Let $\tau$ be the permutation of $\mathbb{Z}/15 = \mathbb{Z}/3 \times \mathbb{Z}/5$ given by $(0)(3)(5)(6)(9)(10)(12)(1\, 7\, 13\, 4\, 11\, 8\, 14\, 2)$. Verify or refute: does $\tau^5$ commute with $\tau^7$ at every point? Identify which CRT face each acts on.

\item \textbf{Run the random control.} Implement the random-control test of Proposition \ref{prop:control}. Report the empirical commutativity probability. Compute the theoretical probability (the number of commuting pairs over the total number of permutation pairs) and compare.

\item \textbf{Find the centralizer.} Within $S_{10}$ (the symmetric group on $\mathbb{Z}/10$), how many permutations $\pi$ commute with $\sigma$? Justify your count by enumerating cycle structures or by computing the orbit-stabilizer relation.

\item \textbf{Different cycle structure.} Replace the $6$-cycle in $\sigma$ with a different one (say, $(0)(3)(8)(9)(1\,4\,2\,7\,5\,6)$). Verify that the new $\sigma^2$ and $\sigma^3$ still commute (powers always do), but examine whether the CRT decomposition still cleanly separates the binary and ternary actions. What is the analogous proposition?

\item \textbf{The role of $\gcd$.} If $\gcd(m, n) > 1$, the CRT bijection (\ref{eq:crt}) fails. Construct a permutation $\sigma$ of $\mathbb{Z}/12 = \mathbb{Z}/4 \times \mathbb{Z}/3$ where the analog of Proposition \ref{prop:proj} either still holds or fails. What changes when the factors are not coprime?

\end{enumerate}

\section{Conclusion}

The Chinese Remainder Theorem is often taught as a computational technique: given two congruences with coprime moduli, find their joint solution. This paper has illustrated a different aspect of the theorem: as a \emph{structural decomposition} that organizes the action of any permutation of $\mathbb{Z}/(mn)$ into factors on $\mathbb{Z}/m \times \mathbb{Z}/n$. The specific example of $\sigma = (0)(3)(8)(9)(1\,7\,6\,5\,4\,2)$ on $\mathbb{Z}/10$ is small enough for a 50-minute class, with a function table that fits on one slide and a verification script that runs in under a second. The commutativity $\sigma^3 \circ \sigma^2 = \sigma^2 \circ \sigma^3$ is trivial as a fact about powers and structurally meaningful via CRT --- a contrast that, for many students, is the first concrete encounter with the idea that \emph{abstract structure} (the product factor) controls \emph{concrete behavior} (the commutativity).

For instructors wanting more examples in this style, the natural next steps are: (a) any small composite modulus $n = pq$ with $p, q$ distinct primes; (b) any cyclic group $\mathbb{Z}/n$ where you can construct a permutation whose order divides $n$ and which has clean CRT-projected behavior; (c) the wreath-product generalizations for more than two prime factors. The pedagogical payoff in every case is the same: a structural decomposition theorem, verifiable computationally before being believed, that organizes apparent coincidence into deductive consequence.

\section*{Acknowledgments}

The author thanks the participants in the CK project for many discussions of substrate-level mathematical structure that surfaced this example.

\appendix

\section{Verification of commutativity at every point}
\label{app:verification}

We compute $\sigma^3(\sigma^2(x))$ and $\sigma^2(\sigma^3(x))$ for each $x \in \mathbb{Z}/10$:

\[
\begin{array}{c|c|c|c}
x & \sigma^2(x) & \sigma^3(\sigma^2(x)) & \sigma^2(\sigma^3(x)) \\
\hline
0 & 0 & 0 & 0 \\
1 & 6 & 2 & 2 \\
2 & 7 & 4 & 4 \\
3 & 3 & 3 & 3 \\
4 & 1 & 5 & 5 \\
5 & 2 & 6 & 6 \\
6 & 4 & 7 & 7 \\
7 & 5 & 1 & 1 \\
8 & 8 & 8 & 8 \\
9 & 9 & 9 & 9
\end{array}
\]

The two columns agree at every $x$, so $\sigma^3 \circ \sigma^2 = \sigma^2 \circ \sigma^3$ as functions.

\section{The verification script}
\label{app:script}

The Python script below performs the verification of Theorem \ref{thm:main} and runs the random control of Proposition \ref{prop:control}.

\begin{verbatim}
# verification_script.py -- runs in under a second
SIGMA = {0:0, 3:3, 8:8, 9:9, 1:7, 7:6, 6:5, 5:4, 4:2, 2:1}

def compose(p, q):
    return {x: p[q[x]] for x in p}

def power(p, n):
    r = {x: x for x in p}
    for _ in range(n):
        r = compose(p, r)
    return r

s2 = power(SIGMA, 2)
s3 = power(SIGMA, 3)
ab = compose(s3, s2)
ba = compose(s2, s3)
assert all(ab[x] == ba[x] for x in SIGMA), "commutativity must hold"
print("Verified: sigma^3 . sigma^2 = sigma^2 . sigma^3 at every point.")
\end{verbatim}

The full script (with the random control) is available as supplementary material and at the project repository.

\begin{thebibliography}{9}

\bibitem{cammann1960}
Cammann, S.\ (1960).
The Evolution of Magic Squares in China.
\emph{Journal of the American Oriental Society}, 80(2), 116--124.
\textit{(historical context for related cyclic-group structure on small finite sets)}

\bibitem{dummit-foote}
Dummit, D.~S., and Foote, R.~M.\ (2004).
\emph{Abstract Algebra}, 3rd Edition.
Wiley.
\textit{(standard reference for the Chinese Remainder Theorem as ring isomorphism --- Chapter 7)}

\bibitem{gauss-dq}
Gauss, C.~F.\ (1801).
\emph{Disquisitiones Arithmeticae}.
\textit{(the modern formulation of CRT as a structural decomposition of $\mathbb{Z}/n$)}

\bibitem{sunzi}
Sun Z\u{\i}\ (3rd century).
\emph{Sun Z\u{\i} Suan Jing} (Master Sun's Mathematical Manual).
\textit{(historical origin of the CRT remainder problem)}

\end{thebibliography}

\end{document}
